% ============================================================
% ETDS Patient Plan Yearly QA Report
% Vorlage - die Platzhalter in doppelten spitzen Klammern setzt create_report.py ein.
% Compiler: XeLaTeX oder LuaLaTeX (wegen fontspec)
% ============================================================
\documentclass[11pt,a4paper]{article}

\usepackage[a4paper,margin=1.8cm,top=2.2cm,bottom=2.2cm]{geometry}
\usepackage{fontspec}
% gleiche Schriftart wie die matplotlib-Plots (qa_plots.py); ohne Arial (z.B. Linux) TeX Gyre Heros
\IfFontExistsTF{Arial}{\setmainfont{Arial}\setsansfont{Arial}}{\setmainfont{TeX Gyre Heros}\setsansfont{TeX Gyre Heros}}
\renewcommand{\familydefault}{\sfdefault}
\usepackage{xcolor}
\usepackage{array}
\usepackage{tabularx}
\usepackage{booktabs}
\usepackage{colortbl}
\usepackage{fancyhdr}
\usepackage{graphicx}
\usepackage{adjustbox}
\usepackage{lastpage}
\usepackage{pdfpages}
% hyperref zuletzt laden (Konvention) - hidelinks: klickbare Links ohne farbige Rahmen,
% damit das Druckbild unveraendert bleibt.
\usepackage[hidelinks]{hyperref}
\usepackage{url}

% ---------- Farben ---------- (barrierefrei)
\definecolor{extdblue}{RGB}{18,52,95}
\definecolor{headerbar}{RGB}{20,60,110}
\definecolor{passgreen}{RGB}{198,224,180}
\definecolor{watchyellow}{RGB}{255,235,156}
\definecolor{actionred}{RGB}{248,203,173}
\definecolor{passgreendark}{RGB}{112,173,71}
\definecolor{tablegray}{RGB}{217,217,217}
\definecolor{extdlabel}{RGB}{31,73,125}

% ---------- Pass/Watch/Action Zellen ----------
\newcommand{\PassCell}{\cellcolor{passgreen}\textbf{Pass}}
\newcommand{\WatchCell}{\cellcolor{watchyellow}\textbf{Watch}}
\newcommand{\ActionCell}{\cellcolor{actionred}\textbf{Action}}

% ---------- Daten aus der Auswertung (Skript-gefuellt) ----------
\newcommand{\Linac}{<<LINAC>>}
\newcommand{\ReportDate}{<<MEAS_DATE>>}          % Messdatum
\newcommand{\ReportTime}{<<MEAS_TIME>>}          % Messbeginn
\newcommand{\CreationDate}{<<CREATION_DATE>>}    % Erstellungsdatum des Reports
\newcommand{\CreationTime}{<<CREATION_TIME>>}
\newcommand{\Passrate}{<<OVERALL_PASSRATE>>}
\newcommand{\ReportAuthor}{<<AUTHOR>>}
\newcommand{\Institution}{<<INSTITUTION>>}
\newcommand{\Phantom}{<<PHANTOM>>}

% Briefkopf-Block: steht im Detailreport rechts neben den Abkuerzungen. create_report.py setzt
% ihn aus Rollen-Code und lokaler Personenliste (interner Report) bzw. dem oeffentlichen
% Kontakt (--public) zusammen; leer, wenn nichts konfiguriert ist.
\newcommand{\ContactBlock}{%
<<CONTACT_BLOCK>>
}

% ---------- Kopf-/Fußzeile ----------
\pagestyle{fancy}
\fancyhf{}
\renewcommand{\headrulewidth}{0pt}
\renewcommand{\footrulewidth}{0.4pt}
\fancyfoot[L]{\footnotesize Report created on \CreationDate\ \CreationTime\\ Report created by \ReportAuthor}
\fancyfoot[R]{\footnotesize <<FOOTER_URL>>\\ Page \thepage\ of \pageref{LastPage}}

% Wiederkehrender Report-Kopf (auf jeder Seite)
\newcommand{\ReportHeader}[1]{%
  \begin{center}
  \begin{tabular}{@{}p{6cm}p{5.5cm}p{4.5cm}@{}}
    {\Huge\bfseries\color{extdblue} ETDS} &
    \begin{tabular}[c]{@{}c@{}}
      \small ETDS QA Report yearly\\
      \small Linac \Linac\ \\
      \small \ReportDate
    \end{tabular} &
    \hfill
\begin{tabular}[l]{@{}c@{}}
  {\bfseries\itshape }\\
  <<LOGO>>
\end{tabular}
  \end{tabular}
  \end{center}
  \vspace{2mm}
  \noindent\colorbox{headerbar}{\makebox[\dimexpr\linewidth-2\fboxsep\relax][l]{\hspace{2mm}\textcolor{white}{\bfseries #1}\hspace{2mm}}}
  \vspace{3mm}
}

% Hilfsmakro: graue Titelzeile über einer Tabelle (einheitlich für ALLE Tabellen)
\newcommand{\SectionBar}[1]{%
  \noindent\colorbox{tablegray}{\makebox[\dimexpr\linewidth-2\fboxsep\relax][c]{\bfseries #1}}\\[1mm]
}

\begin{document}

% ================= SEITE 1: Overview =================
\ReportHeader{ETDS QA Report YEARLY}

% --- QA Details ---
\noindent
\SectionBar{QA Details}
\begin{tabularx}{\linewidth}{@{}l X l >{\raggedleft\arraybackslash}X@{}}
\textcolor{extdlabel}{Tested by} & \ReportAuthor & & \\
\textcolor{extdlabel}{Phantom}   & \Phantom      & \textcolor{extdlabel}{Linac} & \Linac \\
\textcolor{extdlabel}{Date}      & \ReportDate   & \textcolor{extdlabel}{Time}  & \ReportTime \\
\end{tabularx}

\vspace{4mm}

% -------------------------- Ampel-Box + Passraten ------------------------------------------
\SectionBar{Passrates for all 6 DOF}
\noindent
\begin{minipage}[t]{0.2\linewidth}
\centering
\colorbox{<<VERDICT_COLOR>>}{%
  \begin{minipage}{\linewidth}
    \centering
    \small \ReportDate \quad <<MAX_ABS_ERROR>>\\[1mm]
    {\bfseries\Large <<VERDICT_LABEL>>}\\[1mm]
    \small M \quad \Passrate \%
  \end{minipage}%
}
\end{minipage}%
\hspace{0.04\linewidth}%
\begin{minipage}[t]{0.74\linewidth}
\begin{tabularx}{\linewidth}{@{}X r r r c@{}}
\toprule
\textbf{DOF} & \textbf{pass [\%]} & \textbf{watch [\%]} & \textbf{action [\%]} & \textbf{results} \\
\midrule
<<TABLE_PASSRATE_SUMMARY>>
\bottomrule
\end{tabularx}
\end{minipage}
\vspace{5mm}

%-------------------------maximum error metrics ------------------------------
\SectionBar{Maximum error metrics }
\begin{tabularx}{\linewidth}{@{}X X c@{}}
\toprule
\textbf{Max. MAE} & \textbf{Max. RMSE} & \textbf{Max. Absolute Error} \\
\midrule
<<MAX_MAE>> & <<MAX_RMSE>> & <<MAX_ABS_ERROR_UNIT>> \\
\bottomrule
\end{tabularx}
\vspace{5mm}

\SectionBar{Applied tolerance levels and evaluation criteria}
\begin{tabularx}{\linewidth}{@{}X l l l@{}}
\toprule
\textbf{Metric} & \textbf{Pass} & \textbf{Watch} & \textbf{Action} \\
\midrule
\textcolor{extdlabel}{Tracking performance (translation)} &
$\leq <<TOL_ACCEPT>>$ mm &
$<<TOL_ACCEPT>> - <<TOL_WATCH>> $ mm &
$> <<TOL_WATCH>>$ mm \\
\textcolor{extdlabel}{Tracking performance (rotation)} &
$\leq <<TOL_ACCEPT>>^\circ$ &
$<<TOL_ACCEPT>> - <<TOL_WATCH>>  ^\circ$ &
$> <<TOL_WATCH>> ^\circ$ \\
\midrule
\textbf{Criterion} &  &  &  \\
\midrule
\textcolor{extdlabel}{Pass rate [\%]} &
$\geq 99$ &
95--99 &
$<95$ \\
\textcolor{extdlabel}{Watch rate [\%]} &
$\leq 1$ &
1--5 &
-- \\
\textcolor{extdlabel}{Action rate [\%]} &
0 &
0 &
$>0$ \\
\bottomrule
\end{tabularx}

\vspace{5mm}

\noindent
\textbf{Abbrevations:} \\
DOF = degree of freedom \\
MAE = mean average error \\
RMSE = root mean square error \\

%-----------------Unterschrift (nur QMP; leer ohne Freigabe)-----------------
<<SIGNATURE_BLOCK>>

\clearpage

% ================= SEITE 2: Passraten je Messung =================
\ReportHeader{QA Report Details}

\SectionBar{Description }
<<DESCRIPTION>>
\vspace{4mm}

%------------------passrate detailed-----------------------------
\SectionBar{Passrates for all measurments}
\begin{tabularx}{\linewidth}{@{}X X l r r r l@{}}
\toprule
\textbf{DOF} & \textbf{Deflection} & \textbf{Temp.} & \textbf{pass [\%]} & \textbf{watch [\%]} & \textbf{action [\%]} & \textbf{result} \\
\midrule
<<TABLE_PASSRATE_DETAIL>>
\bottomrule
\end{tabularx}
\clearpage

% ================= SEITE 3: Numerische Fehler =================
\ReportHeader{QA Report Details}

\SectionBar{Error Metrics}
\begin{minipage}[t]{\linewidth}
\begin{tabularx}{\linewidth}{@{}X X X X X X r@{}}
\toprule
\textbf{DOF} & \textbf{Unit} & \textbf{Deflection} & \textbf{Temp.} & \textbf{MAE} & \textbf{RMSE} & \textbf{Max Error} \\
\midrule
<<TABLE_ERROR_METRICS>>
\bottomrule
\end{tabularx}
\end{minipage}

%-------------------maximum error metrics compared with previous years ------------------
\vspace{4mm}

\SectionBar{Maximum error metrics compared with previous years }
\begin{tabularx}{\linewidth}{@{}l *{3}{>{\centering\arraybackslash}X}@{}}
\toprule
\textbf{Year} &
\textbf{Max. MAE} &
\textbf{Max. RMSE} &
\textbf{Max. Absolute Error} \\
\midrule
<<TABLE_HISTORY>>
\bottomrule
\end{tabularx}
\vspace{5mm}

\noindent
\textbf{Abbrevations:} \\
DOF = degree of freedom \\
MAE = mean average error \\
RMSE = root mean square error \\

\clearpage

% ================= ANHANG: QA-Uebersichtsplots =================
% Die A4-Plotseiten aus 'etds_qa_evaluation.py plotqa' werden hier 1:1 eingebunden.
% Fuß-/Kopfzeile des Reports bleibt ueber pagecommand erhalten.
<<PLOT_PAGES>>


\end{document}
